\section{自动形式化}

\subsection{两个子问题}

\begin{enumerate}
    \item \textbf{声明形式化}：非形式化定理 $\to$ 形式化声明
    \item \textbf{证明形式化}：非形式化证明 $\to$ 形式化证明
\end{enumerate}

\subsection{声明形式化的评估难题}

同一定理可有多种等价的形式化表述（如"无穷多素数"有多种等价写法），但检验两个形式化声明的逻辑等价性\textbf{在一般情况下是不可判定的}。

\subsection{证明形式化的推理缺口}

人类证明即使很严格，也常有"留给读者"的部分，模型必须自主填补这些\textbf{推理缺口}。

\subsection{案例：欧几里得几何中的自动形式化}

选择欧几里得几何作为受限领域，利用领域特性使原本不可解的问题变得可解。

\subsection{本章小结}

定理证明的核心挑战是如何高效探索无限动作空间。领域特定的洞察可使动作空间结构化，但通用化仍是开放问题。

